% Source: physical pp. 142--153 (printed pp. 119--130). Coverage: \S18.2 from its heading on p. 142 through the two concluding paragraphs at the top of p. 153, ending immediately before \S18.3; Eqs. (18.2.1)--(18.2.50), all unnumbered displays, citations [1,3,3a,4,4a], and two footnotes. Uncertainties: none.
\subsection{The Sliding Scale}

We have seen in the previous section that the large logarithms that appear
at high energy in suitably integrated cross sections or off-shell Feynman
amplitudes can be traced to the prescription used to define renormalized
coupling constants and operators. The central idea of the renormalization
group method is to change this prescription.

Suppose we find some way of defining a new kind of renormalized coupling
constant $g(\mu)$, that depends on a sliding energy scale $\mu$, but that (at
least for $\mu\gg m$) has no dependence on the scale $m$ of the masses of the
theory. Then suitably integrated cross sections or other infrared-safe rate
parameters may be expressed as functions of $g_\mu$ and $\mu$ instead of
$g_R$. By dimensional analysis such functions may be written as
\begin{equation}
\Gamma(E,x,g_\mu,m,\mu)
=E^D\Gamma\left(1,x,g_\mu,\frac{m}{E},\frac{\mu}{E}\right).
\tag{18.2.1}\label{eq:18.2.1}
\end{equation}
\par\noindent
(Our notation here is the same as in Section 18.1; in particular, $x$ stands
for all dimensionless angles, energy ratios, etc. on which $\Gamma$ may
depend.) Since $\mu$ is a completely arbitrary renormalization scale, we can
choose $\mu=E$, in which case Eq.~\eqref{eq:18.2.1} reads
\begin{equation}
\Gamma(E,x,g_\mu,m,\mu)
=E^D\Gamma\left(1,x,g_E,\frac{m}{E},1\right).
\tag{18.2.2}\label{eq:18.2.2}
\end{equation}
\par\noindent
This now has no zero-mass singularities because $g_E$ does not depend on $m$
for $m\ll E$, so there are no large logarithms, and we can use perturbation
theory to calculate $\Gamma$ in terms of $g_E$ as long as $g_E$ itself remains
sufficiently small. In particular, in any finite order of perturbation theory
$\Gamma$ has the asymptotic behavior, for $E\gg m$,
\begin{equation}
\Gamma(E,x,g_\mu,m,\mu)\longrightarrow E^D\Gamma(1,x,g_E,0,1).
\tag{18.2.3}\label{eq:18.2.3}
\end{equation}
\par\noindent
(Non-perturbative corrections are considered in Section 18.4.)

It remains to calculate $g_E$. For instance, in the scalar field theory with
Lagrangian \eqref{eq:18.1.2}, we may define $g_\mu$ in terms of the value of
the scattering amplitude at a renormalization point
$s=t=u=-\mu^2$:
\begin{align}
g_\mu&\equiv A(s=t=u=-\mu^2)\notag\\
&=g-\frac{3g^2}{32\pi^2}\int_0^1dx
\left\{\ln\left(\frac{\Lambda^2}{m^2+\mu^2x(1-x)}\right)-1\right\}
+O(g^3).
\tag{18.2.4}\label{eq:18.2.4}
\end{align}
\par\noindent
or, in terms of the conventional renormalized coupling \eqref{eq:18.1.4},
\begin{equation}
g_\mu=g_R+\frac{3g_R^2}{32\pi^2}\int_0^1dx
\ln\left(1+\frac{\mu^2x(1-x)}{m^2}\right)+O(g_R^3).
\tag{18.2.5}\label{eq:18.2.5}
\end{equation}
\par\noindent
But this formula is reliable only if the correction term is smaller than
$g_R$; that is, only if $\lvert g_R\ln(\mu/m)\rvert\ll1$. If this were the
case for $\mu\simeq E$, then we would not need the methods of the
renormalization group; ordinary perturbation theory would be good enough.

Instead of using formulas like \eqref{eq:18.2.5} directly for large $\mu$, we
must instead proceed in stages: $g_\mu$ may be calculated in terms of $g_R$
as long as $\mu/m$ is not much larger than unity; then $g_{\mu'}$ may be
calculated in terms of $g_\mu$ as long as $\mu'/\mu$ is not much larger than
unity; and so on, up to $g_E$. Instead of discrete stages, this may also be
done continuously. Dimensional analysis tells us that the relation between
$g_{\mu'}$ and $g_\mu$ takes the form
\begin{equation}
g_{\mu'}=G(g_\mu,\mu'/\mu,m/\mu).
\tag{18.2.6}\label{eq:18.2.6}
\end{equation}
\par\noindent
Differentiating with respect to $\mu'$ and then setting $\mu'=\mu$ yields the
differential equation
\begin{equation}
\mu\frac{d}{d\mu}g_\mu=\beta\left(g_\mu,\frac{m}{\mu}\right),
\tag{18.2.7}\label{eq:18.2.7}
\end{equation}
\par\noindent
where
\begin{equation}
\beta\left(g_\mu,\frac{m}{\mu}\right)
\equiv\left[\frac{\partial}{\partial z}G(g_\mu,z,m/\mu)\right]_{z=1}.
\tag{18.2.8}\label{eq:18.2.8}
\end{equation}
\par\noindent
There are no zero-mass singularities here, so for $\mu\gg m$ the differential
equation becomes simply
\begin{equation}
\mu\frac{d}{d\mu}g_\mu=\beta(g_\mu,0)\equiv\beta(g_\mu),
\tag{18.2.9}\label{eq:18.2.9}
\end{equation}
\par\noindent
which is often known as the Callan--Symanzik equation~\hyperref[ch18-ref-1]{[1]}.
We are to calculate $g_E$ by integrating the differential equation
\eqref{eq:18.2.9}, with an initial value $g_M$ at some scale $\mu=M$, chosen
in practice large enough so that for $\mu\geq M$ we can neglect the masses
$m$ compared with $\mu$, but small enough so that large logarithms
$\ln(M/m)$ do not prevent us from using perturbation theory to calculate
$g_M$ in terms of the conventional renormalized coupling constant $g_R$. The
solution may be formally written
\begin{equation}
\ln(E/M)=\int_{g_M}^{g_E}\frac{dg}{\beta(g)}
\tag{18.2.10}\label{eq:18.2.10}
\end{equation}
\par\noindent
as long as $\beta(g)$ does not vanish between $g_M$ and $g_E$.

The results of the previous paragraph do not rely on perturbation theory, but
we usually need to use perturbation theory to calculate the functions $G$ and
$\beta$. As an example, suppose we calculate $g_{\mu'}$ in the scalar field
theory with interaction $g\phi^4/24$, renormalizing by expressing $g$ in terms
of $g_\mu$ rather than $g_R$. Following the same procedure that led to
Eq.~\eqref{eq:18.2.5}, this gives
\[
g_{\mu'}=g_\mu-\frac{3g_\mu^2}{32\pi^2}\int_0^1dx
\ln\left(\frac{m^2+\mu^2x(1-x)}{m^2+\mu'^2x(1-x)}\right)
+O(g_\mu^3).
\]
Then Eq.~\eqref{eq:18.2.8} gives
\begin{equation}
\beta\left(g_\mu,\frac{m}{\mu}\right)
=+\frac{3g_\mu^2}{16\pi^2}\int_0^1dx
\frac{\mu^2x(1-x)}{m^2+\mu^2x(1-x)}+O(g_\mu^3).
\tag{18.2.11}\label{eq:18.2.11}
\end{equation}
\par\noindent
For $\mu\gg m$, this is
\begin{equation}
\beta(g_\mu)=\frac{3g_\mu^2}{16\pi^2}+O(g_\mu^3).
\tag{18.2.12}\label{eq:18.2.12}
\end{equation}
\par\noindent
In next order, the beta-function for $\mu\gg m$ is~\hyperref[ch18-ref-3]{[3]}
\[
\beta(g_\mu)=g_\mu\left[3\left(\frac{g_\mu}{16\pi^2}\right)
-\frac{17}{3}\left(\frac{g_\mu}{16\pi^2}\right)^2+\cdots\right].
\]

If we are content with the one-loop approximation the calculation of
$\beta(g)$ can be done even more easily. In order to avoid large radiative
corrections in matrix elements at energies of order $\mu$, we must write the
bare coupling $g$ in terms of a finite renormalized coupling $g_\mu$, as
\begin{equation}
g=g_\mu+B(g_\mu)\ln\frac{\Lambda}{\mu}+\cdots.
\tag{18.2.13}\label{eq:18.2.13}
\end{equation}
\par\noindent
For instance, from Eq.~\eqref{eq:18.1.3} we could have immediately read off
that the $\ln\Lambda$ terms in $g$ have coefficient
\begin{equation}
B(g)=-\frac{3}{2}\left[-i(2\pi)^4g^2\right]
\left\{\left[\frac{-i}{(2\pi)^4}\right]^2 2\pi^2i\right\}
=\frac{3g^2}{16\pi^2}.
\tag{18.2.14}\label{eq:18.2.14}
\end{equation}
\par\noindent
The unrenormalized coupling is of course independent of $\mu$, so to lowest
order
\[
\mu\frac{d}{d\mu}g_\mu-B(g_\mu)=0
\]
and so to lowest order
\begin{equation}
\beta(g)=B(g).
\tag{18.2.15}\label{eq:18.2.15}
\end{equation}
\par\noindent
With $B(g)$ given by Eq.~\eqref{eq:18.2.14}, this agrees with our previous
result \eqref{eq:18.2.12} for $\beta(g)$ in the $\phi^4$ scalar field theory.

Instead of using a simple ultraviolet cut-off, which can interfere with gauge
invariance, it is often more convenient to deal with ultraviolet divergences
by use of dimensional regularization. For a spacetime dimensionality $d<4$,
we find in place of $\ln(\Lambda/\mu)$ the convergent integral
\[
\int_\mu^\infty k^{d-4}\frac{dk}{k}
=\frac{\mu^{d-4}}{4-d}
\underset{d\to4}{\longrightarrow}
\left[\frac{1}{4-d}-\ln\mu\right].
\]
Thus instead of eliminating the cut-off dependence by writing the
unrenormalized coupling constant as in \eqref{eq:18.2.13}, we instead write
\begin{equation}
g=g_\mu+B(g_\mu)\left[\frac{1}{4-d}-\ln\mu\right]
\tag{18.2.16}\label{eq:18.2.16}
\end{equation}
\par\noindent
with the same function $B(g_\mu)$ as before. Thus in order to calculate
$\beta(g_\mu)$, all we need to do is to pick out the coefficient of the
singular factor $1/(4-d)$ in the renormalized coupling. This argument is
extended to all orders of perturbation theory in Section 18.6.

As long as $g_\mu$ is sufficiently small in the scalar field theory with
Lagrangian \eqref{eq:18.1.2}, the solution of Eqs.~\eqref{eq:18.2.9} and
\eqref{eq:18.2.12} can be well approximated by
\begin{equation}
g_\mu=-\frac{16\pi^2}{3\ln(\mu/M)},
\tag{18.2.17}\label{eq:18.2.17}
\end{equation}
\par\noindent
where $M$ is an integration constant. This expression illustrates a common
aspect of renormalization group calculations, that dimensionless couplings
like $g_R$ become replaced with parameters like $M$ that have the
dimensionality of mass. The value of $M$ may be related to $g_R$ by comparing
the solution \eqref{eq:18.2.17} with the behavior of the coupling for values
of $\mu$ that are large enough to allow us to use approximations based on
$\mu\gg m$, but small enough so that
$\lvert g_R\ln(\mu/m)\rvert\ll1$, where \eqref{eq:18.2.5} gives
\begin{equation}
g_\mu\simeq g_R+\frac{3g_R^2}{16\pi^2}\ln\left(\frac{\mu}{m}\right).
\tag{18.2.18}\label{eq:18.2.18}
\end{equation}
\par\noindent
In this way, we find
\begin{equation}
M\simeq m\exp\left(\frac{16\pi^2}{3g_R}\right),
\tag{18.2.19}\label{eq:18.2.19}
\end{equation}
\par\noindent
so that Eq.~\eqref{eq:18.2.17} may be put in a more conventional form
\begin{equation}
g_\mu=g_R\left[1-\frac{3g_R}{16\pi^2}\ln\frac{\mu}{m}\right]^{-1}.
\tag{18.2.20}\label{eq:18.2.20}
\end{equation}
\par\noindent
To repeat, this is valid provided $g_\mu$ is small, even if
$\lvert g_R\ln(\mu/m)\rvert$ is of order unity, so it represents a significant
improvement over the perturbative result \eqref{eq:18.2.18}. Of course, the
condition that $g_\mu$ should be small will become violated when
$g_R\ln(\mu/m)$ is sufficiently close to the critical value $16\pi^2/3$. But
at least Eq.~\eqref{eq:18.2.20} makes the unequivocal prediction that $g_E$
becomes large enough to invalidate perturbation theory at some energy $E$
below the critical value \eqref{eq:18.2.19}.

In calculating off-shell matrix elements of operators instead of integrated
cross sections, we also need to take into account the $N$-factors that appear
in the definition of the renormalized operators whose matrix elements are
finite. We saw in the previous section that if these $N$-factors are defined
in a conventional way (say, so that the correction factors produced by the
divergent subgraphs are cancelled when the operator carries zero
four-momentum, or is a field on its mass shell) then the formula for the
$N$-factor involves zero-mass singularities as in Eq.~\eqref{eq:18.1.12} or
Eq.~\eqref{eq:18.1.16}, resulting in large logarithms at energies $E\gg m$.
The cure is to define renormalization constants $N_\mu^{(\mathcal O)}$ at a
sliding scale $\mu$, so that in matrix elements of the renormalized operator
\begin{equation}
\mathcal O_\mu=N_\mu^{(\mathcal O)}\mathcal O
\tag{18.2.21}\label{eq:18.2.21}
\end{equation}
\par\noindent
the correction factor produced by divergent subgraphs containing operator
$\mathcal O$ are cancelled at a renormalization point characterized by
four-momenta of order $\mu$. If $M_R$ is a matrix element of operators that
are conventionally renormalized, and $M$ is one in which the operators are
renormalized as in Eq.~\eqref{eq:18.2.21}, then for any $\mu$
\begin{equation}
M_R=\left[\prod_{\mathcal O}
\left(N^{(\mathcal O)}/N_\mu^{(\mathcal O)}\right)\right]
M(E,x,g_\mu,m,\mu).
\tag{18.2.22}\label{eq:18.2.22}
\end{equation}
\par\noindent
We can again use dimensional analysis (assuming $M$ has dimensionality $D$)
and set $\mu=E$, to write this as
\begin{equation}
M_R=E^D\left[\prod_{\mathcal O}
\left(N^{(\mathcal O)}/N_E^{(\mathcal O)}\right)\right]
M\left(1,x,g_E,\frac{m}{E},1\right).
\tag{18.2.23}\label{eq:18.2.23}
\end{equation}
\par\noindent
Thus to find the high energy behavior of an off-shell amplitude $M_R$, we need
to know how $N_\mu$ varies with the renormalization scale $\mu$.

For any two renormalization scales $\mu$ and $\mu'$, the renormalized
operators $N_\mu^{(\mathcal O)}\mathcal O$ and
$N_{\mu'}^{(\mathcal O)}\mathcal O$ both have finite matrix elements, so the
ratio $N_{\mu'}^{(\mathcal O)}/N_\mu^{(\mathcal O)}$ must be cut-off
independent. On dimensional grounds, this ratio must take the form
\begin{equation}
N_{\mu'}^{(\mathcal O)}/N_\mu^{(\mathcal O)}
=G^{(\mathcal O)}(g_\mu,\mu'/\mu,m/\mu).
\tag{18.2.24}\label{eq:18.2.24}
\end{equation}
\par\noindent
Differentiating with respect to $\mu'$ and then setting $\mu'=\mu$ gives
\begin{equation}
\mu\frac{d}{d\mu}N_\mu^{(\mathcal O)}
=\gamma^{(\mathcal O)}(g_\mu,m/\mu)N_\mu^{(\mathcal O)},
\tag{18.2.25}\label{eq:18.2.25}
\end{equation}
\par\noindent
where
\begin{equation}
\gamma^{(\mathcal O)}(g_\mu,m/\mu)
\equiv\left[\frac{\partial}{\partial z}
G^{(\mathcal O)}(g_\mu,z,m/\mu)\right]_{z=1}.
\tag{18.2.26}\label{eq:18.2.26}
\end{equation}
\par\noindent
The solution is
\begin{equation}
N_E^{(\mathcal O)}\propto
\exp\left[\int^E\gamma^{(\mathcal O)}
\left(g_\mu,\frac{m}{\mu}\right)\frac{d\mu}{\mu}\right].
\tag{18.2.27}\label{eq:18.2.27}
\end{equation}
\par\noindent
This is a useful result because the introduction of the sliding scale
prevents the appearance of zero-mass singularities in
$N_\mu^{(\mathcal O)}$ and $N_{\mu'}^{(\mathcal O)}$, and hence also in
$G^{(\mathcal O)}$ and $\gamma^{(\mathcal O)}$. Hence as long as $g_\mu$ is
small, there are no large logarithms that prevent the application of
perturbation theory to calculate $\gamma^{(\mathcal O)}$. Also, for
$\mu\gg m$, $\gamma^{(\mathcal O)}(g_\mu,m/\mu)$ has a smooth limit
\begin{equation}
\gamma^{(\mathcal O)}(g_\mu)
\equiv\gamma^{(\mathcal O)}(g_\mu,0).
\tag{18.2.28}\label{eq:18.2.28}
\end{equation}
\par\noindent

As an example, consider the operator $\mathcal O=\phi^2$ in the scalar field
theory with interaction $g\phi^4/24$. Instead of renormalizing it so that the
correction factor \eqref{eq:18.1.9} is cancelled at $p^2=0$, we cancel it at a
sliding scale $p^2=\mu^2$, by introducing a new renormalized $\phi^2$ operator
$N_\mu^{(\phi^2)}\phi^2$, with
\[
N_\mu^{(\phi^2)}\equiv F^{(\phi^2)}(\mu^2)^{-1}
=1+\frac{g}{32\pi^2}\int_0^1dx
\left[\ln\left(\frac{\Lambda^2}{m^2+\mu^2x(1-x)}\right)-1\right]
+O(g^2).
\]
Then the function \eqref{eq:18.2.24} for this operator is
\[
\begin{aligned}
G^{(\phi^2)}(g_\mu,\mu'/\mu,m/\mu)
&\equiv\frac{N_{\mu'}^{(\phi^2)}}{N_\mu^{(\phi^2)}}\\
&=1+\frac{g_\mu}{32\pi^2}\int_0^1dx
\ln\left[\frac{m^2+\mu^2x(1-x)}{m^2+\mu'^2x(1-x)}\right]
+O(g_\mu^2).
\end{aligned}
\]
(We can use $g_\mu$ here instead of $g$ or $g_{\mu'}$, because the difference
only affects the terms of order $g_\mu^2$ or higher.) From
Eq.~\eqref{eq:18.2.26}, we have then
\[
\gamma^{(\phi^2)}\left(g_\mu,\frac{m}{\mu}\right)
=-\frac{g_\mu}{16\pi^2}\int_0^1
\frac{\mu^2x(1-x)}{m^2+\mu^2x(1-x)}\,dx+O(g_\mu^2)
\]
or for $\mu\gg m$,
\begin{equation}
\gamma^{(\phi^2)}(g_\mu)=-\frac{g_\mu}{16\pi^2}+O(g_\mu^2).
\tag{18.2.29}\label{eq:18.2.29}
\end{equation}
\par\noindent

Another good example is provided by the $N$-factor associated with the
renormalization of the electromagnetic field in quantum electrodynamics.
Recall that the photon propagator can be made finite for all momenta by
evaluating it for renormalized electromagnetic fields, or equivalently by
multiplying the propagator of the unrenormalized field by $Z_3^{-1}$:
\begin{equation}
\widetilde\Delta_{\rho\sigma}(q)=Z_3^{-1}\Delta_{\rho\sigma}(q).
\tag{18.2.30}\label{eq:18.2.30}
\end{equation}
\par\noindent
Eq.~(10.5.17) shows that this renormalized propagator may be written
\begin{equation}
\widetilde\Delta_{\rho\sigma}(q)
=\frac{\eta_{\rho\sigma}}{[q^2-i\epsilon][1-\pi(q^2)]}
+q_\rho q_\sigma\text{-terms}.
\tag{18.2.31}\label{eq:18.2.31}
\end{equation}
\par\noindent
Suppose we instead define a renormalized field $N_\mu^{(A)}A^\rho$, whose
propagator has a term proportional to $\eta_{\rho\sigma}/[q^2-i\epsilon]$
with a coefficient that is equal to unity at a sliding renormalization scale
$q^2=\mu^2$. For this purpose, we must clearly take
\begin{equation}
N_\mu^{(A)}=Z_3^{-1/2}[1-\pi(\mu^2)]^{1/2}.
\tag{18.2.32}\label{eq:18.2.32}
\end{equation}
\par\noindent
Using Eq.~(11.2.22), the function \eqref{eq:18.2.24} is then
\begin{align}
G^{(A)}(g_\mu,\mu'/\mu,m/\mu)
&=\left[\frac{1-\pi(\mu'^2)}{1-\pi(\mu^2)}\right]^{1/2}\notag\\
&=1-\frac{e_\mu^2}{4\pi^2}\int_0^1dx\,x(1-x)
\ln\left[\frac{m^2+\mu'^2x(1-x)}{m^2+\mu^2x(1-x)}\right]
+O(e_\mu^4)
\tag{18.2.33}\label{eq:18.2.33}
\end{align}
\par\noindent
and so Eq.~\eqref{eq:18.2.26} gives
\begin{equation}
\gamma^{(A)}(e_\mu,m/\mu)
=-\frac{e_\mu^2}{2\pi^2}\int_0^1dx\,
\frac{x^2(1-x)^2\mu^2}{m^2+\mu^2x(1-x)}+O(e_\mu^4).
\tag{18.2.34}\label{eq:18.2.34}
\end{equation}
\par\noindent
As promised, this has a smooth limit for $\mu\gg m$
\begin{equation}
\gamma^{(A)}(e_\mu)\equiv\gamma^{(A)}(e_\mu,0)
=-\frac{e_\mu^2}{12\pi^2}+O(e_\mu^4).
\tag{18.2.35}\label{eq:18.2.35}
\end{equation}
\par\noindent

As already mentioned, electrodynamics is a special case, because the
renormalization constant $Z_3^{-1/2}$ in the definition of the renormalized
electromagnetic field is just the reciprocal of the constant used to define
the renormalized electric charge of the electron: $e_R=Z_3^{1/2}e$. The
natural definition of the renormalized electric charge at a sliding scale
$\mu$ is then
\begin{equation}
e_\mu=N_\mu^{(A)-1}e=Z_3^{-1/2}N_\mu^{(A)-1}e_R,
\tag{18.2.36}\label{eq:18.2.36}
\end{equation}
\par\noindent
\begin{sloppypar}
so that $e_\mu$ times the field $N_\mu^{(A)}A^\rho$ renormalized at scale
$\mu$ is independent of $\mu$. From Eq.~\eqref{eq:18.2.25} we see that the
function $\beta(e)$ which gives the $\mu$ dependence of $e_\mu$ according to
Eq.~\eqref{eq:18.2.9} for $\mu\gg m$ has the value
\end{sloppypar}
\begin{equation}
\beta(e)=-e\gamma^{(A)}(e)=\frac{e^3}{12\pi^2}+O(e^5).
\tag{18.2.37}\label{eq:18.2.37}
\end{equation}
\par\noindent
Using an earlier calculation~\hyperref[ch18-ref-3a]{[3a]} of the fourth-order
term in the vacuum polarization function $\pi(q^2)$, Gell-Mann and Low were
able also to give the term in $\beta(e)$ of next order in $e$:
\begin{equation}
\beta(e)=\frac{e^3}{12\pi^2}+\frac{e^5}{64\pi^4}+O(e^7).
\tag{18.2.38}\label{eq:18.2.38}
\end{equation}
\par\noindent
In other words, the electric charge at a sliding scale $\mu$ satisfies the
renormalization group equation
\begin{equation}
\mu\frac{d}{d\mu}e_\mu
=\frac{e_\mu^3}{12\pi^2}+\frac{e_\mu^5}{64\pi^4}+O(e_\mu^7).
\tag{18.2.39}\label{eq:18.2.39}
\end{equation}
\par\noindent
This shows, for small $e_\mu$, that $e_\mu$ \emph{increases} with increasing
$\mu$.

We also need an initial condition. This is provided by the known value of the
conventionally renormalized charge $e_R=Z_3^{1/2}e$, for which
$\alpha\equiv e_R^2/4\pi=1/137.036\ldots$. Eqs.~\eqref{eq:18.2.32} and
\eqref{eq:18.2.36} yield
\begin{align}
e_R/e_\mu
&=Z_3^{1/2}N_\mu^{(A)}=\sqrt{1-\pi(\mu^2)}\notag\\
&=1-\frac{e_R^2}{4\pi^2}\int_0^1dx\,x(1-x)
\ln\left[1+\frac{\mu^2x(1-x)}{m_e^2}\right]+O(e_R^4).
\tag{18.2.40}\label{eq:18.2.40}
\end{align}
\par\noindent
We need to match this with the solution of Eq.~\eqref{eq:18.2.39} at a value
of $\mu$ which is large enough to justify the approximation $\mu\gg m_e$ in
\eqref{eq:18.2.39} but small enough so that the logarithm in
Eq.~\eqref{eq:18.2.40} is still small enough compared with $4\pi^2/e_R^2$ to
justify the use of perturbation theory. (For instance, we might take $\mu$ to
be of order 100 MeV.) For such values of $\mu$, Eq.~\eqref{eq:18.2.40} gives
\begin{equation}
e_\mu\simeq e_R+\frac{e_R^3}{12\pi^2}
\left[\ln\frac{\mu}{m_e}-\frac{5}{6}\right].
\tag{18.2.41}\label{eq:18.2.41}
\end{equation}
\par\noindent
On the other hand, the solution of Eq.~\eqref{eq:18.2.39} for $e_\mu$ small
(keeping only the leading term on the right-hand side) is
\begin{equation}
e_\mu=\left[\text{constant}-\frac{\ln\mu}{6\pi^2}\right]^{-1/2}.
\tag{18.2.42}\label{eq:18.2.42}
\end{equation}
\par\noindent
Comparing Eqs.~\eqref{eq:18.2.41} and \eqref{eq:18.2.42} gives the solution
\begin{equation}
e_\mu=e_R\left[1-\frac{e_R^2}{6\pi^2}
\left(\ln\left(\frac{\mu}{m_e}\right)-\frac{5}{6}\right)\right]^{-1/2}.
\tag{18.2.43}\label{eq:18.2.43}
\end{equation}
\par\noindent
Unlike Eq.~\eqref{eq:18.2.41}, Eq.~\eqref{eq:18.2.43} is valid as long as
$e_\mu^2/6\pi^2$ is small, whether or not
$(e_R^2/6\pi^2)\allowbreak\ln(\mu/m_e)$ is small.

For instance, we have already seen in Section 11.3 that the leading
fourth-order radiative correction to the magnetic moment of the muon can be
obtained by multiplying the second-order (Schwinger) term (11.3.16)
by the vacuum polarization function $\pi_e(k^2)$ at $k^2\simeq m_\mu^2$,
which according to Eq.~\eqref{eq:18.2.40} is the same (to this order) as using
$e_{m_\mu}^2/4\pi$ in place of $\alpha$ in the Schwinger term.

Another example: Experiments at high energy electron--positron colliders such
as LEP at CERN or SLC at SLAC now study physical processes at energies of the
order of the mass of the $Z^0$ particle, or 91 GeV. Eq.~\eqref{eq:18.2.43}
shows that at these energies, radiative corrections in pure quantum
electrodynamics should be calculated using a value for the fine structure
constant which is not $\alpha=1/137.036$ but rather
\begin{equation}
\frac{e^2(91\ \mathrm{GeV})}{4\pi}
=\frac{\alpha}{1-2(11.25)\alpha/3\pi}=\frac{1}{134.6}.
\tag{18.2.44}\label{eq:18.2.44}
\end{equation}
\par\noindent
This is for a theory in which electrons are the only charged particles with
masses below $m_Z$. In the real world there are many such particle types, and
the effective fine structure constant~\hyperref[ch18-ref-4]{[4]} at $m_Z$ is
$(128.87\pm0.12)^{-1}$.

\begin{center}
$*\ *\ *$
\end{center}

The sliding scale at which the coupling parameters of a theory are calculated
may be the value of an external \emph{field} rather than the momentum of an
external line. In one of the early applications of the renormalization group
method, Coleman and E.~Weinberg~\hyperref[ch18-ref-4a]{[4a]} considered the
effective potential $V(\phi)$ for a spacetime-independent external scalar
field $\phi$. In the simple case where the field interacts with itself alone,
their one-loop result is given by Eq.~\eqref{eq:16.2.15}. It is particularly
interesting to consider this potential in the case where the renormalized
mass $m_R$ vanishes, where to one-loop order we may set $\mu^2(\phi)$ in the
last term equal to $g_R\phi^2/2$, so that Eq.~\eqref{eq:16.2.15} here reads
(with $g$ a slightly redefined coupling):
\begin{equation}
V(\phi)=\lambda_R+\frac{g}{24}\phi^4
+\frac{g^2\phi^4\ln\phi^2}{256\pi^2}.
\tag{18.2.45}\label{eq:18.2.45}
\end{equation}
\par\noindent
This looks like at first sight as if for $g>0$ the potential becomes less than
$\lambda_R$ for very small $\phi$, so that the point $\phi=0$ is a local
maximum instead of a minimum, but for such small values of $\phi$ the third
term is larger than the second, and the perturbation theory is obviously
untrustworthy. Also, we would like to be able to argue that $g$ must be
positive in order for the potential to be bounded below at large fields, but
Eq.~\eqref{eq:18.2.45} shows that however small $g$ is, there is some
sufficiently large $\phi$ where perturbation theory breaks down, and hence
Eq.~\eqref{eq:18.2.45} cannot be trusted to tell us whether the potential is
bounded below at large fields.

We can do much better by using a coupling constant defined at a sliding scale
$\mu$ of field strength. Suppose we define a coupling $g_\mu$ by the condition
that
\begin{equation}
V(\mu)=\lambda_R+\frac{g_\mu}{24}\mu^4.
\tag{18.2.46}\label{eq:18.2.46}
\end{equation}
\par\noindent
If we had used $g_\mu$ as the coupling parameter from the beginning, then in
place of Eq.~\eqref{eq:18.2.45} we would have obtained
\begin{equation}
V(\phi)=\lambda_R+\frac{g_\mu}{24}\phi^4
+\frac{g_\mu^2\phi^4}{256\pi^2}\ln\left(\frac{\phi^2}{\mu^2}\right),
\tag{18.2.47}\label{eq:18.2.47}
\end{equation}
\par\noindent
which obviously satisfies Eq.~\eqref{eq:18.2.46}.
\footnote{Eq.~\eqref{eq:18.2.47} differs from what we would get from
Eq.~\eqref{eq:18.2.45} by simply using Eq.~\eqref{eq:18.2.46} to express $g$
in terms of $g_\mu$, in that the last term is proportional to $g_\mu^2$ rather
than $g^2$. The difference is of higher order in $g$, but can become
significant if $V(\phi)$ is evaluated at $\phi$ very different from $\mu$,
where large logarithms can compensate for powers of coupling. If we use
$g_\mu$ as the coupling parameter from the beginning and take $\mu$ to be of
order $\phi$ then no such large logarithms occur, and the approximation
Eq.~\eqref{eq:18.2.47} is valid as long as $g_\mu$ remains small.}
The renormalization group equation for $g_\mu$ can be obtained from the
condition that this effective potential is independent of $\mu$,
\footnote{To eliminate all cut-off dependence, $\phi$ should be written in
terms of a renormalized field $\phi_\mu=N_\mu^\phi\phi$, which gives
$V(\phi_\mu)$ a $\mu$ dependence arising from the $\mu$ dependence of the
renormalization constant $N_\mu^\phi$. This point is ignored here, because in
the scalar field theory with interaction $\propto\phi^4$ the
lowest-order graphs contributing to the $\mu$ dependence of $N_\mu^\phi$ have
two loops, so that to the order of the calculations presented here we can
take $N_\mu^\phi=1$.}
\begin{equation}
\mu\frac{dg_\mu}{d\mu}=\frac{3g_\mu^2}{16\pi^2}.
\tag{18.2.48}\label{eq:18.2.48}
\end{equation}
\par\noindent
(Terms involving the derivative of $g_\mu^2$ are dropped here, because they
are of higher order in $g_\mu$, and are hence negligible as long as $g_\mu$
is sufficiently small.) It is not a coincidence that this takes the same
form as the renormalization group equation \eqref{eq:18.2.9},
\eqref{eq:18.2.12}, where $\mu$ was a renormalization momentum, because as we
shall see in the next section the first two terms in the renormalization group
equation are always independent of the way that we define the sliding scale.
The solution of this equation is given by Eq.~\eqref{eq:18.2.17}, in general
with a different integration constant $M$. Hence by taking $\mu=\phi$ in
Eq.~\eqref{eq:18.2.17}, we now have
\begin{equation}
V(\phi)=\lambda_R-\frac{32\pi^2\phi^4}{3\ln(\phi^2/M^2)}.
\tag{18.2.49}\label{eq:18.2.49}
\end{equation}
\par\noindent
\begin{equation}
g_\phi=-\frac{32\pi^2}{3\ln(\phi^2/M^2)}.
\tag{18.2.50}\label{eq:18.2.50}
\end{equation}
\par\noindent
This result should be used with some care, because it is only valid where the
coupling constant $g_\phi$ is small. The problem is not just that
Eq.~\eqref{eq:18.2.49} loses its validity for $\phi$ near $M$; we also cannot
integrate the renormalization group equation through the singularity at
$\phi=M$, so a knowledge of $g_\phi$ on one side of this singularity tells us
nothing about its behavior on the other side.

If $g_{\phi_0}$ is found to have a small \emph{positive} value for some
$\phi_0$ then from Eq.~\eqref{eq:18.2.50} we know that $M>\lvert\phi_0\rvert$,
so $g_\phi$ remains small and Eq.~\eqref{eq:18.2.49} is valid for
$\lvert\phi\rvert<\lvert\phi_0\rvert$. This shows that the point $\phi=0$ is
a local \emph{minimum} of $V(\phi)$, contrary to what Eq.~\eqref{eq:18.2.45}
might have led us to suppose. This means that the vacuum which is invariant
under the symmetry transformation $\phi\to-\phi$ is stable in this model,
aside from the possibility of quantum mechanical barrier penetration. On the
other hand, we do not know that Eq.~\eqref{eq:18.2.49} becomes valid for
$\lvert\phi\rvert$ sufficiently large compared with $M$. The coupling might
remain too large to use perturbation theory for all $\lvert\phi\rvert>M$, and
even if it becomes small for some $\lvert\phi\rvert>M$, the potential might
be given for such $\phi$ by Eq.~\eqref{eq:18.2.49} with a renormalization scale
$M'>\phi$, producing a second singularity. Thus in this case we cannot
conclude that $V(\phi)\to-\infty$ for $\lvert\phi\rvert\to\infty$.

Similarly, if $g_\phi$ is found to have a small \emph{negative} value for some
$\phi_0$ then we know that $M<\lvert\phi_0\rvert$, so $g_\phi$ remains small
and Eq.~\eqref{eq:18.2.49} is valid for
$\lvert\phi\rvert>\lvert\phi_0\rvert$. In this case we cannot conclude
anything about the behavior of the potential for $\lvert\phi\rvert<M$, but
here we \emph{can} use Eq.~\eqref{eq:18.2.49} to see that
$V(\phi)\to-\infty$ for $\lvert\phi\rvert\to\infty$, ruling out the
possibility of any stable vacuum. Because we are here considering the limit
$\lvert\phi\rvert\to\infty$, this conclusion holds also for scalar fields of
mass $m>0$, provided $m\ll\lvert\phi_0\rvert$. This is why it is necessary to
assume that the $\phi^4$ coupling (renormalized at any scale much larger than
the scalar mass) is positive.
